% Preface -- Financial Engineering and the Evidence
\chapter*{Preface}
\addcontentsline{toc}{chapter}{Preface}
\markboth{Preface}{Preface}

\emph{Financial Engineering and the Evidence: Active Learning by Verifying Famous Claims}

\section*{A natural way to learn}

For a long time I have been thinking about how people learn a technical subject, and how that learning could be more active and more inspiring.

Many textbooks follow a familiar pattern: a model is presented, its formulas are derived, a few examples are worked, and exercises follow at the end of the chapter. This is a time-tested way to pass on knowledge.

But it may not make a reader research-oriented, and it may not build the judgment a working quant needs most: knowing whether a result can be trusted. When I turn the question around and ask, ``Does this actually work? Let's find out,'' the classroom comes alive. Students make predictions, argue about what the data will show, run the numbers, and remember what they found, often because it was not what they expected.

I believe this is the natural way people learn, even in the age of AI. This book is built on that conviction. Its subtitle, \emph{Active Learning by Verifying Famous Claims}, describes the method. You will not be asked to accept the results of quantitative finance on authority. Instead, you will be invited to test those results for yourself while learning the techniques of financial engineering.

\section*{Why this approach, and why now}

Three developments make this the right moment for a book like this.

First, the evidence behind many famous claims is live and unresolved. For example, Hou, Xue and Zhang (2020) re-tested 452 published anomalies and found that about 65 percent fail a basic significance test once micro-caps are given little weight. Jensen, Kelly and Pedersen (2023), using a Bayesian method and data from 93 countries, concluded that most factors do replicate. Goyal, Welch and Zafirov (2024) found that many return predictors proposed since 2008 have lost their power. When leading researchers disagree, students need to be able to work through the disagreement themselves. This book is designed to bring students into such debates rather than only report their conclusions. In Chapter~\ref{ch:07}, for example, students replicate both sides of the anomaly debate and form their own view.

Second, the data to do so is now free. Kenneth French's data library, FRED, the Jensen--Kelly--Pedersen global factors and Chen and Zimmermann's Open Source Asset Pricing portfolios let any student rebuild published results without a paid database.

Third, new methods are entering finance faster than they are being tested. Machine learning and time series foundation models promise much, but the finance-specific evidence is mixed. Students need to learn where these tools help and where they do not, and the only way to know is to test them.

Courses that teach the replication debate usually do it through journal articles, while their textbooks present results rather than ask students to re-test them. This book is written to fill that gap.

\section*{Putting famous claims on trial}

Finance is full of famous claims. Past winners keep winning. The trend is your friend. Stock market returns are predictable. Optimized portfolios beat simple equal weighting. Machine learning predicts stock returns better than linear models.

Practitioners act on these claims, products are sold on them, and many are taught as if they were settled. They are also exactly the kind of question an active learner can answer with data.

In this book, every such claim is put on trial, and you are the one who tries it. Each claim chapter follows the same cycle:

\begin{enumerate}
  \item \textbf{Claim.} State what is asserted, in one testable sentence.
  \item \textbf{Understand.} Learn how the idea works and how the original evidence was produced.
  \item \textbf{Verify.} Rebuild the result in a Python lab on open data, then stress-test it out of sample, in other markets and after costs.
  \item \textbf{Judge.} Reach a dated verdict: survives, weakens, fails or provisional.
  \item \textbf{Apply.} Take away what investors and practitioners can use, and what remains open.
\end{enumerate}

This order is the heart of the active approach. Because the question comes first, every method is learned at the moment it is needed, to settle something you actually want to know. Newey--West errors, false-discovery rates and walk-forward evaluation stop being abstract tools and become the instruments that decide a case. Each lab also asks you to write down what you expect before you run a step. When the data disagree with your prediction, the surprise is visible, and it stays with you.

The verdicts are not foregone conclusions. Some claims largely survive, such as the predictability of volatility. Some survive only in weakened form, such as value and momentum. The goal is to follow the evidence wherever it leads, and on some questions, such as whether most published anomalies are false, serious researchers still disagree. There you will replicate both sides and write your own verdict.

Repeated across twenty-three claims, the cycle becomes a habit. That habit, more than any single verdict, is what I hope you take from this book: before believing a result, ask ``compared to what, out of sample, and after costs?''

\section*{How the book is organized}

The book has eight parts, twenty-three chapters and four appendices. Every chapter, from the first, puts a famous claim on trial, and the chapters are ordered so that each one introduces methods only a little harder than the ones before it. Read in order, the book is a single course that starts with averages and t-statistics and ends with machine learning and foundation models.

\begin{itemize}
  \item Part I tests claims about returns, luck and data, using means, Sharpe ratios, t-statistics and simulation.
  \item Part II tests claims about the cross-section of stock returns, using regression, factor models, portfolio sorts and trading costs.
  \item Part III tests claims about forecasting over time, in stocks, bonds and currencies, using out-of-sample evaluation.
  \item Part IV tests claims about portfolio construction and diversification.
  \item Part V tests claims about volatility and risk models.
  \item Part VI tests claims about options and credit.
  \item Part VII tests claims about machine learning, from tree ensembles to transformer-based asset-pricing models.
  \item Part VIII tests claims about time series foundation models and ends with a capstone in which you replicate a claim of your own choosing.
\end{itemize}

The parts are not marked by separate pages in the text. Table~\ref{tab:pref-parts} lists the chapters in each part.

\begin{table}[htbp]
\centering
\small
\caption{The eight parts of the book and their chapters.}
\label{tab:pref-parts}
\begin{tabularx}{0.95\linewidth}{@{}>{\raggedright\arraybackslash}p{0.07\linewidth}>{\raggedright\arraybackslash}p{0.28\linewidth}>{\raggedleft\arraybackslash}p{0.04\linewidth}>{\raggedright\arraybackslash}X@{}}
\toprule
Part & Theme & Ch. & Title \\
\midrule
I & Returns, luck and data & \ref{ch:01} & A Strong Track Record Proves Skill \\
 &  & \ref{ch:02} & Historical Data Show What a Strategy Would Have Earned \\
\addlinespace
II & Factor models and the cross-section & \ref{ch:03} & Size and Value Explain Returns That Beta Cannot \\
 &  & \ref{ch:04} & A Profitable Backtest Means a Profitable Strategy \\
 &  & \ref{ch:05} & Low-Risk Stocks Beat High-Risk Stocks \\
 &  & \ref{ch:06} & Past Winners Keep Winning \\
 &  & \ref{ch:07} & Most Published Anomalies Fail to Replicate \\
\addlinespace
III & Forecasting over time & \ref{ch:08} & Stock Market Returns Are Predictable \\
 &  & \ref{ch:09} & The Trend Is Your Friend \\
 &  & \ref{ch:10} & The Yield Curve Predicts Future Interest Rates \\
 &  & \ref{ch:11} & Currency Losses Cancel Out Higher Interest Rates \\
\addlinespace
IV & Portfolios and dependence & \ref{ch:12} & Optimized Portfolios Beat Simple Equal Weighting \\
 &  & \ref{ch:13} & Diversification Fails When You Need It Most \\
\addlinespace
V & Volatility and risk & \ref{ch:14} & Volatility Is Predictable \\
 &  & \ref{ch:15} & Value-at-Risk Models Measure Risk Reliably \\
\addlinespace
VI & Derivatives and credit & \ref{ch:16} & Selling Options Earns a Reliable Premium \\
 &  & \ref{ch:17} & Delta Hedging Removes Option Risk \\
 &  & \ref{ch:18} & Corporate Bond Spreads Mainly Reflect Expected Defaults \\
\addlinespace
VII & Machine learning & \ref{ch:19} & Machine Learning Predicts Stock Returns Better Than Linear Models \\
 &  & \ref{ch:20} & Artificial Intelligence Prices Stocks Better Than Factor Models \\
\addlinespace
VIII & Foundation models and the capstone & \ref{ch:21} & Foundation Models Forecast Better Than Classical Methods \\
 &  & \ref{ch:22} & Foundation Models Trained on Market Data Beat GARCH \\
 &  & \ref{ch:23} & Capstone: Replicating a Finance Claim \\
\bottomrule
\end{tabularx}
\end{table}


The claims were chosen so that testing them carries you through the standard curriculum of financial engineering. Each one rests on a core model: the CAPM, GARCH, Black--Scholes, the yield curve, the Merton model, mean--variance optimization. Every claim chapter opens with a self-contained Foundations section that teaches that model before the claim is tested. You meet the same material as in a conventional course, but you meet each model at the point where it makes a prediction the data can contradict.

This also means the book can stand on its own, or sit beside the classics. Hull, Cochrane, and Campbell, Lo and MacKinlay teach how the models are built. This book asks whether their predictions hold, and the table later in this preface shows, for each chapter, where those texts treat the topic in depth.

Every chapter has a Python lab. The book describes each lab by its question, data, setup and results, with tables and charts, so you can follow every argument without running any code. The code itself lives in a companion GitHub repository, where it can be kept working as data sources and libraries change. All labs run on free data, including the Kenneth French library, FRED and the open replication datasets now available, so no paid subscription is required.

\section*{Why these claims}

Each of the twenty-three claims, in Chapters~\ref{ch:01} to~\ref{ch:22}, was chosen to do two jobs at once: to settle a question that matters, and to teach a core model. The table shows how the claims map onto the standard curriculum. For each area of financial engineering, it lists the model or tool that the chapter's Foundations section teaches and the claim that puts it to work.

\begingroup
\small
\setlength{\LTpre}{6pt}\setlength{\LTpost}{6pt}
\begin{longtable}{@{}>{\raggedright\arraybackslash}p{0.198\linewidth}>{\raggedright\arraybackslash}p{0.288\linewidth}>{\raggedright\arraybackslash}p{0.297\linewidth}r@{}}
\toprule
Core area & Model or tool & Claim on trial & Ch. \\
\midrule
\endfirsthead
\toprule
Core area & Model or tool & Claim on trial & Ch. \\
\midrule
\endhead
\bottomrule
\endlastfoot
Statistics of returns & Returns, Sharpe ratio, t-statistics, luck by simulation & A strong track record proves skill & \ref{ch:01} \\
Financial data & Adjusted prices, survivorship, look-ahead and micro-cap biases & Historical data show what a strategy would have earned & \ref{ch:02} \\
Asset pricing and factors & OLS, the CAPM, factor models, robust standard errors & Size and value explain returns that beta cannot & \ref{ch:03} \\
Trading and implementation & Spreads, turnover, break-even costs, capacity & A profitable backtest means a profitable strategy & \ref{ch:04} \\
CAPM and leverage & Fama--MacBeth, GRS, beta estimation, leverage constraints & Low-risk stocks beat high-risk stocks & \ref{ch:05} \\
Market efficiency and behavior & Portfolio sorts, cross-sectional momentum, crash risk & Past winners keep winning & \ref{ch:06} \\
Econometrics of finance & False discovery rate, shrinkage, post-publication tests & Most published anomalies fail to replicate & \ref{ch:07} \\
Forecasting & Walk-forward evaluation, out-of-sample $R^2$, forecast comparison tests & Stock market returns are predictable & \ref{ch:08} \\
Technical analysis & Trading rules, time-series momentum, the Reality Check and SPA tests & The trend is your friend & \ref{ch:09} \\
Fixed income & Present value, yield curves, forward rates, duration, expectations hypothesis & The yield curve predicts future interest rates & \ref{ch:10} \\
International finance & Covered and uncovered interest parity, carry & Currency losses cancel out higher interest rates & \ref{ch:11} \\
Portfolio construction & Mean--variance, efficient frontier, covariance estimation & Optimized portfolios beat simple equal weighting & \ref{ch:12} \\
Diversification and dependence & Portfolio variance, correlation dynamics, tail dependence & Diversification fails when you need it most & \ref{ch:13} \\
Volatility modelling & ARCH/GARCH, realized and implied volatility & Volatility is predictable & \ref{ch:14} \\
Risk management & VaR, expected shortfall, backtesting, Basel & VaR models measure risk reliably & \ref{ch:15} \\
Options & Payoffs, put--call parity, implied volatility & Selling options earns a reliable premium & \ref{ch:16} \\
Derivative pricing and hedging & Black--Scholes, the Greeks, dynamic hedging & Delta hedging removes option risk & \ref{ch:17} \\
Credit & Default, recovery, spreads, the Merton model & Corporate bond spreads mainly reflect expected defaults & \ref{ch:18} \\
Machine learning & Regularization, trees, neural nets, cross-validation & Machine learning predicts stock returns better than linear models & \ref{ch:19} \\
Asset pricing with AI & Stochastic discount factor, IPCA, attention and transformers & Artificial intelligence prices stocks better than factor models & \ref{ch:20} \\
Forecasting with foundation models & ARIMA, exponential smoothing, zero-shot transformers & Foundation models forecast better than classical methods & \ref{ch:21} \\
Financial foundation models & Fine-tuning, range-based volatility, synthetic market data & Foundation models trained on market data beat GARCH & \ref{ch:22} \\
\end{longtable}
\endgroup

Each method is taught in the first chapter that needs it and is reused in the chapters that follow.

Every claim had to pass five tests before it earned a chapter.

\begin{enumerate}
  \item \textbf{Anchored in a core model.} Testing it requires you to learn a model that belongs in any financial engineering course.
  \item \textbf{Famous and consequential.} Practitioners act on it, products are sold on it, or it shaped the research literature.
  \item \textbf{Testable with free data.} Any student can run the lab with the French library, FRED, Shiller's data and similar public sources.
  \item \textbf{Still debated.} Recent papers disagree, so the verdict is not a foregone conclusion and your own test adds information.
  \item \textbf{A spread of verdicts.} Some claims largely survive (volatility is predictable), some survive only in weakened form (value, momentum, the option premium), and some largely fail (forward rates as forecasts, optimized portfolios beating an equal-weighted one). The method does not always say no.
\end{enumerate}

The claims do not cover everything. This is not a full treatment of stochastic calculus, exotic derivatives, term-structure models beyond the basics, corporate finance, or market microstructure at the level of the order book. 

\section*{Foundations sections by chapter}

Every chapter except the capstone opens its teaching with a Foundations section. It is not optional background: it sets out the concepts, formulas and terms that the rest of the chapter depends on, and it is part of what each chapter expects you to learn. Each Foundations section runs about two to three pages, with key terms in bold, the essential formulas and a worked example. The table below lists, for each chapter, the standard texts that treat its topic in depth. They are written so that a reader without a prior course in investments or financial engineering can follow every chapter, which lets the book stand on its own. You do not need the books in the Further reading column to follow this one. They are listed for readers who want a deeper treatment of a topic, now or later.

\begingroup
\small
\setlength{\LTpre}{6pt}\setlength{\LTpost}{6pt}
\begin{longtable}{@{}r>{\raggedright\arraybackslash}p{0.468\linewidth}>{\raggedright\arraybackslash}p{0.324\linewidth}@{}}
\toprule
Ch. & Foundations covers & Further reading \\
\midrule
\endfirsthead
\toprule
Ch. & Foundations covers & Further reading \\
\midrule
\endhead
\bottomrule
\endlastfoot
\ref{ch:01} & Simple, log and excess returns; the risk-free rate; expected return and volatility; the Sharpe ratio; standard errors and t-statistics & Bodie, Kane and Marcus (2024); Campbell, Lo and MacKinlay (1997) \\
\ref{ch:02} & Corporate actions and adjusted prices; index weighting; CRSP, Compustat, FRED and the French library; delisting; point-in-time data & Bodie, Kane and Marcus (2024); Bali, Engle and Murray (2016) \\
\ref{ch:03} & Moments of a distribution; OLS and its standard errors; the CAPM, beta and alpha; factor models & Wooldridge (2025); Bodie, Kane and Marcus (2024); Cochrane (2005) \\
\ref{ch:04} & Order books; bid--ask spreads; market and limit orders; short selling; long--short portfolios; turnover & Harris (2003); Bodie, Kane and Marcus (2024); Hull (2022) \\
\ref{ch:05} & The security market line; estimating beta; Fama--MacBeth and GRS tests; leverage constraints & Bodie, Kane and Marcus (2024); Ang (2014) \\
\ref{ch:06} & Measuring past returns; cross-sectional momentum; underreaction and overreaction & Bodie, Kane and Marcus (2024); Ilmanen (2011) \\
\ref{ch:07} & Market efficiency; what an anomaly is; replication and robustness; false discovery rates & Bodie, Kane and Marcus (2024); Bali, Engle and Murray (2016) \\
\ref{ch:08} & Loss functions; information sets; bias--variance; training, validation and test data; autoregressive models; walk-forward evaluation and leakage; out-of-sample $R^2$ and forecast comparison tests & Hyndman and Athanasopoulos (2021); L\'opez de Prado (2018) \\
\ref{ch:09} & Moving-average and breakout rules; time-series momentum; volatility targeting; the SPA test & Ilmanen (2011); Elliott and Timmermann (2016) \\
\ref{ch:10} & Present value; bond prices and yields; zero-coupon curves; forward rates; duration; the expectations hypothesis & Hull (2022); Tuckman and Serrat (2022) \\
\ref{ch:11} & Spot and forward exchange rates; covered and uncovered interest parity; carry & Hull (2022); Bodie, Kane and Marcus (2024) \\
\ref{ch:12} & Mean--variance optimization; the efficient frontier; estimating covariance matrices & Bodie, Kane and Marcus (2024); Luenberger (2014) \\
\ref{ch:13} & Portfolio variance and correlation; correlation over time; tail dependence & Bodie, Kane and Marcus (2024); Ang (2014) \\
\ref{ch:14} & Realized and implied volatility; ARCH and GARCH; the VIX & Hull (2022); Tsay (2010) \\
\ref{ch:15} & Value-at-risk and expected shortfall; historical, parametric and Monte Carlo methods; backtesting; Basel rules & Hull (2022, 2023) \\
\ref{ch:16} & Calls and puts; payoffs; put--call parity; implied volatility; the variance risk premium & Hull (2022) \\
\ref{ch:17} & The Black--Scholes model and its assumptions; delta, gamma and the Greeks; dynamic hedging & Hull (2022); Shreve (2004) \\
\ref{ch:18} & Default probability and recovery; credit spreads; ratings; the Merton model & Hull (2022) \\
\ref{ch:19} & Supervised learning; preparing inputs; penalized regression; trees, random forests and boosting; neural networks; tuning and interpreting models & James, Witten, Hastie and Tibshirani (2021); Hastie, Tibshirani and Friedman (2009); L\'opez de Prado (2018) \\
\ref{ch:20} & The stochastic discount factor; shrinkage and IPCA estimators; deep and adversarial SDFs; attention and transformers & Cochrane (2005); Kelly and Xiu (2023); Nagel (2021) \\
\ref{ch:21} & Classical forecasting models (ARIMA, exponential smoothing); transformers; zero-shot forecasting & Hyndman and Athanasopoulos (2021) \\
\ref{ch:22} & OHLCV data and range-based volatility; fine-tuning; synthetic market data; the Giacomini--White test & Hyndman and Athanasopoulos (2021); Goodfellow, Bengio and Courville (2016) \\
\end{longtable}
\endgroup

Chapter~\ref{ch:23}, the capstone, draws on all of these and has no Foundations section of its own. It also collects the research-design tools (hypotheses, power, pre-registration and robustness grids) that a replication needs.


Full references for these books are listed alphabetically by author under Further Reading at the end of the book (page~\pageref{furtherreading}).


\section*{Who this book is for}

This book is for anyone who builds, buys or studies quantitative strategies and needs to judge whether the evidence behind them holds up.

\begin{itemize}
  \item Students in financial engineering, quantitative finance and financial data science will use the whole book, labs included, as will quantitative analysts and researchers at asset managers, banks and fintech firms.
  \item PhD students can use it as a companion for the replication debate. Portfolio managers, allocators and serious individual investors can read the principles and verdicts and skim the labs.
\end{itemize}

You will need to possess basic Python, and to know introductory statistics, regression and the elements of investments. Everything else is taught along the way.

This is not a book of trading signals. If anything, its most consistent lesson is how rarely a promising backtest survives honest testing. What it offers instead is more durable: the ability to tell the difference.

\section*{An invitation}

Learning by verifying is slower at first than learning by reading. You will make predictions that turn out wrong, and you will spend time on results that do not hold. But you will understand each method because you needed it, and you will remember each verdict because you reached it yourself.

I wrote this book in the hope that it makes learning financial engineering more active and more enjoyable. Pick a claim, write down what you expect, and let the evidence decide.

\bigskip
\begin{flushright}
\emph{Chris Kuo}\\
\emph{January 2027}
\end{flushright}
